\begin{abstract}
The same Python program can be compiled by \tcompile along several paths. Graph breaks are placed at positions
that depend on the call site, automatic differentiation is traced only when a gradient is required, loops are
vectorized only above a length threshold, execution may take place under a dispatch mode, autocast, or CUDA
graphs, and a program may be exported and packaged for deployment. Existing testing techniques for AI compiler
stacks vary the program and compile it along whichever path the harness takes, and they compare output
values. Consistency failures between eager and compiled execution, however, often occur where a rarely used
program feature meets a rarely taken path, and many of them do not change an output value.
This paper varies the compilation instead and treats its path as a test dimension. We identify four path factors that a harness can
enumerate and control one at a time: the position of graph breaks, the backend layer, the execution context,
and the artifact route from export to a packaged runtime. Together with program-side factors derived from the
source and from operator metadata, they form an execution matrix in which every difference from the reference
execution has one candidate cause. Executions are compared on seven observables (values against a
high-precision reference, metadata, object identity, exception type and routing, side effects, gradients, and
process survival), and layered execution names the compiler stage at which a difference first appears.
Applied to PyTorch~2.14 and nightly builds, the approach has produced 73 reports: 55 new issues, 15
issue comments, and 3 reports to a C++ compiler vendor. At the time of writing, 14 are fixed or have a fix under
review, 10 are confirmed, 48 await triage, and 1 was rejected. On a 185-program corpus that program-only
differential testing under three configurations had already covered, one path factor exposed two new defects; one
program-side factor exposed 134 operators that compute under \tcompile on dtypes their eager kernels reject.
Factor families for which no inconsistency was found are reported as well.
\end{abstract}
